
\documentclass[11pt,a4paper]{article}
\usepackage[a4paper,margin=2.05cm]{geometry}
\usepackage[spanish,es-noquoting]{babel}
\usepackage[T1]{fontenc}
\usepackage[utf8]{inputenc}
\usepackage{lmodern,microtype}
\usepackage{amsmath,amssymb,amsthm,mathtools}
\usepackage{booktabs,longtable,array,enumitem}
\usepackage{xcolor,hyperref}
\pagecolor{white}\color{black}
\setlength{\parindent}{0pt}
\setlength{\parskip}{0.65em}
\title{\bfseries N74 --- Monodromía de nueve puertas\\
\large La raíz vuelve; el borde recuerda}
\author{HMT / auditoría técnica}
\date{}
\begin{document}
\maketitle

\begin{abstract}
N74 prueba la lectura sugerida por la novena puerta: la torre no es una sucesión plana ni una periodicidad ordinaria. Es una monodromía de nueve fases. La fase \(g(K)=K\bmod9\) retorna cada nueve puertas, pero el estado de frontera no retorna idéntico: vuelve cargado por carry, cola y supervivencia. La ley visible es: \(\varphi\) fija el eje, \(\pi/e\) forman el espejo móvil, y el futuro selecciona la orientación. La novena puerta no crea la ley; la hace visible.
\end{abstract}

\section{Ley de fase}

Definimos:
\[
g(K)=K\bmod9,\qquad 0\equiv9.
\]
Entonces:
\[
\boxed{
g(K+9)=g(K).
}
\]
Pero la frontera no es periódica:
\[
\boxed{
\mathcal S_{K+9}\neq \mathcal S_K\quad\text{en general}.
}
\]
Así, la torre es:
\[
\boxed{
\text{monodromía, no periodicidad.}
}
\]

\section{Perfil por fase}

\[
\begin{array}{c|c|c|c|c}
\text{fase}&\#&t&\#\text{orient.}&h_{\max}\\
\hline
1 & 1 & 10 & 1 & 2 \\
2 & 3 & 11,20,21 & 3 & 3 \\
3 & 1 & 12 & 1 & 2 \\
4 & 2 & 13,23 & 2 & 3 \\
5 & 3 & 5,14,24 & 3 & 2 \\
6 & 3 & 6,15,25 & 3 & 3 \\
7 & 2 & 7,16 & 2 & 3 \\
8 & 2 & 8,17 & 2 & 3 \\
9 & 1 & 9 & 1 & 2 \\
\end{array}
\]

\section{Retorno \(t\to t+9\)}

\[
\begin{array}{c|c|c|c|c|c|c}
t&t+9&K\bmod9&\text{tipo}_t&\text{tipo}_{t+9}&\varphi=&\pi+e=\\
\hline
5 & 14 & 5 & orientation_gate & orientation_gate & False & False \\
6 & 15 & 6 & orientation_gate & orientation_gate & False & False \\
7 & 16 & 7 & orientation_gate & orientation_gate & False & False \\
8 & 17 & 8 & orientation_gate & orientation_gate & False & False \\
11 & 20 & 2 & orientation_gate & orientation_gate & False & False \\
12 & 21 & 3 & orientation_gate & orientation_gate & False & False \\
14 & 23 & 5 & orientation_gate & orientation_gate & False & False \\
15 & 24 & 6 & orientation_gate & orientation_gate & False & False \\
16 & 25 & 7 & orientation_gate & orientation_gate & False & False \\
\end{array}
\]

La fase retorna, pero los valores de frontera no se repiten. Eso es memoria de carry.

\section{Puerta de orientación}

De N73:
\[
q=u^\pi+u^e+u^\varphi,\qquad
a=u^\pi+u^e-u^\varphi,
\]
por tanto:
\[
\boxed{
u^\varphi=2(q-a),
\qquad
u^\pi+u^e=q-u^\varphi.
}
\]
Si hay ambigüedad local, sólo puede vivir en el reparto \(\pi/e\). Por eso:
\[
\boxed{
\varphi=\text{eje fijo},
\qquad
\pi/e=\text{espejo móvil}.
}
\]

\section{Dictamen}

\[
\boxed{
\text{verde: la ley correcta es de nueve fases con monodromía de carry.}
}
\]
\[
\boxed{
\text{verde: la novena puerta es una puerta de orientación visible, no un stutter.}
}
\]
\[
\boxed{
\text{verde: la raíz vuelve y el borde recuerda.}
}
\]
\[
\boxed{
\text{amarillo: convertir esta monodromía finita auditada en teorema infinito.}
}
\]
\end{document}
